\hspace{3mm}

\centerline{\bf \Huge  Дополнительные построения. Треугольник.}
\centerline{\bf \Huge Уровень: Light.}

\begin{flushright}
    \scriptsize{lightweight babyyyyy}
\end{flushright}


Топ 3 дополнительных построения в треугольнике (без окружностей):

\textit{3 место.} Возьмём прямоугольный треугольник с углами 30, 60, 90. Отразим его относительно большего катета. Полученный треугольник будет правильный.

\textit{2 место.} \textbf{Двойной прогиб медианой.} При удвоении медианы в треугольнике образуется параллелограм.

\textit{1 место.} Из середины стороны проведём прямую параллельную какой-нибудь другой стороне. Получим среднюю линию, которая разбивает наш треугольник на трапецию и ещё один треугольник подобный исходному, с коэфф-ом подобия $\dfrac{1}{2}$

\problem{1}
Докажите, что середины двух противоположных сторон любого четырёхугольника и середины его диагоналей являются вершинами параллелограмма.


\problem{2} 
В треугольнике $ABC$ проведены медиана $BM$ и высота $AH$. Известно, что  $BM = AH$.  Найдите угол $MBC$.

\problem{3} 
На медиане $AA_1$ треугольника $ABC$ взята точка $M$, причём  $AM : MA_1 = 1:3$.  В каком отношении прямая BM делит сторону $AC$?

\problem{4}
В равнобедренном треугольнике с боковой стороной, равной 4, проведена медиана к боковой стороне. Найдите основание треугольника, если медиана равна 3.

\problem{5}
Докажите, что если в треугольнике медиана и биссектриса совпадают, то тре- угольник равнобедренный (используя удвоение медианы).

\problem{6}  На медиане $BM$ треугольника $ABC$ возьмём точку $E$ так, что угол $CEM$ равен углу $ABM$. Тогда отрезок $CE$ равен одной из сторон треугольника.

\problem{7} В треугольнике $ABC$ провели медиану $BM$. Оказалось, что сумма углов $A$ и $C$ равна углу $ABM$. Найдите отношение медианы $BM$ к стороне $BC$.

\problem{8} В треугольнике $ABC$ точка $M$ – середина стороны $BC$, точка $E$ лежит внутри стороны $AC,$  $BE\geq2AM.$  Докажите, что треугольник $ABC$ тупоугольный.
